\documentclass[a4paper,11pt]{article}
\usepackage[T2A]{fontenc}
\usepackage[utf8]{inputenc}
\usepackage[russian]{babel}
\usepackage{amsmath,amssymb,mathtools}
\usepackage{geometry}
\geometry{left=26mm,right=26mm,top=23mm,bottom=23mm}
\usepackage{microtype}
\usepackage{tikz}
\usetikzlibrary{calc,arrows.meta,positioning,shapes.geometric}
\usepackage{tkz-euclide}
\usepackage{booktabs,tabularx,array}
\usepackage{enumitem}
\usepackage{fancyhdr}
\usepackage{hyperref}
\usepackage{parskip}
\usepackage{xcolor}
\usepackage{setspace}
\usepackage{needspace}
\definecolor{accent}{HTML}{008080}
\definecolor{darkaccent}{HTML}{005B5B}
\definecolor{textgray}{HTML}{383838}
\definecolor{quiet}{HTML}{687376}
\definecolor{gold}{HTML}{A36A36}
\hypersetup{colorlinks=true,linkcolor=darkaccent,pdfauthor={Лёвин Артём Александрович},pdftitle={Чек-лист. Прямоугольный параллелепипед и куб}}
\setlength{\parindent}{0pt}
\setlength{\parskip}{6pt}
\setstretch{1.12}
\setlist[itemize]{leftmargin=6mm,itemsep=2pt,topsep=2pt}
\setlist[enumerate]{leftmargin=7mm,itemsep=4pt,topsep=2pt}
\pagestyle{fancy}
\fancyhf{}
\fancyhead[L]{\footnotesize\color{textgray} Прямоугольный параллелепипед и куб}
\fancyhead[R]{\footnotesize\color{textgray} Ксения Клыкова}
\renewcommand{\headrulewidth}{0.35pt}
\renewcommand{\headrule}{\hbox to\headwidth{\color{accent}\leaders\hrule height \headrulewidth\hfill}}
\fancyfoot[C]{\small\thepage}
\setlength{\headheight}{22pt}
\newcommand{\sectionhead}[1]{\Needspace{7\baselineskip}\section{#1}\vspace{-1mm}}
\newcommand{\checkline}[1]{\par\noindent$\square$\hspace{2mm}#1\par}
\newcommand{\keyidea}[1]{\textcolor{darkaccent}{\textbf{Ключевая идея.}} #1}
\newcommand{\note}[1]{\textcolor{gold}{\textbf{Обратите внимание.}} #1}
\newcommand{\sol}[1]{\textcolor{darkaccent}{\textbf{Решение.}} #1}
\newcommand{\ex}[1]{\par\medskip\textcolor{darkaccent}{\textbf{Пример.}} #1\par}
\renewcommand{\thesection}{\arabic{section}}
% No dependency on titlesec: use the standard section formatter.
\makeatletter
\renewcommand\section{\@startsection{section}{1}{0pt}{2.2ex plus .7ex minus .2ex}{1ex plus .2ex}{\normalfont\large\color{darkaccent}}}
\makeatother
\begin{document}
\begin{titlepage}
\thispagestyle{empty}
\begin{center}
\vspace*{20mm}
{\large\color{darkaccent} МАТЕМАТИКА}\par
\vspace{24mm}
{\huge Чек-лист}\par
\vspace{12mm}
{\LARGE\color{accent} Прямоугольный параллелепипед\\[3mm] и куб}\par
\vspace{14mm}
{\normalsize Подготовка к ЕГЭ по профильной математике}\par
\vfill
{\large Ксения Клыкова}\par
\vspace{11mm}
{19 ноября 2026 г.}\par
\vspace{25mm}
{\small Лёвин Артём Александрович}\\[3mm]
{\small эксклюзивно для Ксении Клыковой}\par
\vspace{18mm}
\begin{tikzpicture}[x=1cm,y=1cm]
\draw[accent, line width=.75pt] (0,0) .. controls (1.7,.8) and (3.2,-.8) .. (4.7,0)
 .. controls (6.2,.8) and (7.7,-.8) .. (9.2,0);
\end{tikzpicture}
\vspace{6mm}
\end{center}
\end{titlepage}
\setcounter{page}{1}

\sectionhead{Что изображено на чертеже}
\checkline{Различаю \textbf{ребро}, \textbf{грань}, \textbf{диагональ грани} и \textbf{пространственную диагональ}.}

\begin{center}
\begin{tikzpicture}[scale=0.92]
  \tkzDefPoint(0,0){A}
  \tkzDefPoint(4.4,0){B}
  \tkzDefPoint(5.8,1.3){C}
  \tkzDefPoint(1.4,1.3){D}
  \tkzDefPoint(0,2.8){Aone}
  \tkzDefPoint(4.4,2.8){Bone}
  \tkzDefPoint(5.8,4.1){Cone}
  \tkzDefPoint(1.4,4.1){Done}
  \tkzDrawSegments[dashed,thin,color=quiet](A,D D,C D,Done)
  \tkzDrawSegments[thin](A,B B,C C,Cone Cone,Done Done,Aone Aone,A B,Bone Bone,Cone Aone,Bone)
  \tkzDrawSegment[dashed,accent,line width=1.3pt](A,C)
  \tkzDrawSegment[dashed,gold,line width=1.3pt](A,Cone)
  \tkzDrawPoints[size=2.3,color=darkaccent](A,B,C,D,Aone,Bone,Cone,Done)
  \tkzLabelPoint[below left](A){$A$}
  \tkzLabelPoint[below](B){$B$}
  \tkzLabelPoint[right](C){$C$}
  \tkzLabelPoint[left](D){$D$}
  \tkzLabelPoint[left](Aone){$A_1$}
  \tkzLabelPoint[right](Bone){$B_1$}
  \tkzLabelPoint[above right](Cone){$C_1$}
  \tkzLabelPoint[above](Done){$D_1$}
  \node[quiet,font=\small,anchor=north] at (2.2,-0.22) {$a$};
  \node[quiet,font=\small,anchor=west] at (5.18,.55) {$b$};
  \node[quiet,font=\small,anchor=east] at (-.08,1.4) {$c$};
\end{tikzpicture}

{\footnotesize \textcolor{accent}{\rule{7mm}{1pt}} диагональ грани $AC$ \qquad
\textcolor{gold}{\rule{7mm}{1pt}} пространственная диагональ $AC_1$}
\end{center}

Прямоугольный параллелепипед имеет шесть прямоугольных граней. Его три измерения обозначим $a$, $b$, $c$; все они \textbf{положительны}. На чертеже $AB=a$, $BC=b$, $AA_1=c$.

\note{Пространственная диагональ соединяет противоположные вершины тела, например $A$ и $C_1$. Диагональ грани лежит в одном прямоугольнике, например $A$ и $C$.}

\sectionhead{Объём прямоугольного параллелепипеда}
\checkline{Умею выбирать основание и соответствующую высоту.}
\[
\boxed{V=S_{\text{осн}}\cdot h=abc}
\]
Если основание --- грань $ABCD$, то $S_{\text{осн}}=ab$, а высота равна $c$. При выборе другой грани основанием меняется и соответствующая высота; результат $abc$ остаётся тем же.

\ex{Объём равен $24\ \text{см}^3$, ребро, перпендикулярное искомой грани, равно $3\ \text{см}$. Найдите площадь этой грани.}
\[
V=S\cdot h \quad\Longrightarrow\quad S=\frac{V}{h}=\frac{24}{3}=8\ \text{см}^2.
\]
\keyidea{В текстовой задаче сначала запишите формулу для величины, которую требуется найти; затем определите известные и неизвестные данные.}

\sectionhead{Площадь полной поверхности}
\checkline{Отличаю площадь одной грани от площади всей поверхности.}
\[
\boxed{S_{\text{полн}}=2(ab+ac+bc)}
\]
Каждый из трёх видов прямоугольных граней встречается дважды. Поэтому складываем площади всех шести граней.

\ex{При $a=2$, $b=3$, $c=4$ площадь основания $ab=6$, а площадь всей поверхности}
\[
S_{\text{полн}}=2(6+8+12)=52.
\]
\Needspace{3\baselineskip}\note{В условии слова «площадь грани», «площадь основания» и «площадь полной поверхности» обозначают разные величины.}

\sectionhead{Две разные диагонали}
\checkline{Применяю теорему Пифагора к диагонали грани.}
\[
\boxed{d_{ab}=\sqrt{a^2+b^2}},\qquad
 d_{ac}=\sqrt{a^2+c^2},\qquad d_{bc}=\sqrt{b^2+c^2}.
\]

\checkline{Умею находить пространственную диагональ.}
\[
\boxed{D=\sqrt{a^2+b^2+c^2}},\qquad D^2=a^2+b^2+c^2.
\]
Это следствие \textbf{двукратного} применения теоремы Пифагора: сначала $AC^2=a^2+b^2$, затем $AC_1^2=AC^2+c^2$.

\ex{Два ребра равны $2$ и $4$, пространственная диагональ равна $6$. Найдите объём.}
\[
\begin{aligned}
6^2&=2^2+4^2+c^2,\\
c^2&=36-4-16=16,\qquad c=4,\\
V&=2\cdot4\cdot4=32.
\end{aligned}
\]
\note{Длина ребра положительна, поэтому из $c^2=16$ берём $c=4$. Если в условии указана просто «диагональ параллелепипеда», обычно имеется в виду пространственная диагональ; диагональ конкретной грани называют отдельно.}

\sectionhead{Куб как частный случай}
\checkline{Использую равенство всех двенадцати рёбер куба: $a=b=c$.}
\begin{center}
\renewcommand{\arraystretch}{1.35}
\begin{tabularx}{\linewidth}{@{}X>{\centering\arraybackslash}p{.37\linewidth}@{}}
\toprule
Величина & Формула \\
\midrule
Объём & $V=a^3$ \\
Площадь полной поверхности & $S=6a^2$ \\
Пространственная диагональ & $D=a\sqrt{3}$ \\
Диагональ грани & $d=a\sqrt{2}$ \\
\bottomrule
\end{tabularx}
\end{center}
Формулы для диагоналей следуют из равенств $D^2=3a^2$ и $d^2=2a^2$.

\checkline{Умею находить площадь поверхности \textbf{без вычисления самого ребра}.}
\[
\boxed{S=2D^2}\qquad\text{или}\qquad\boxed{S=3d^2}.
\]
\ex{Если пространственная диагональ куба равна $1$, то $S=2\cdot1^2=2$. Если диагональ грани равна $3$, то $S=3\cdot3^2=27$.}

\sectionhead{Как переводить условие в уравнение}
\checkline{Ввожу обозначения для старого и нового ребра и объёма.}
Если ребро куба \textbf{увеличили на $1$}, то
\[
a_2=a_1+1,\qquad V_1=a_1^3,\qquad V_2=(a_1+1)^3.
\]
Если объём \textbf{увеличился на $19$}, то $V_2-V_1=19$, значит,
\[
(a_1+1)^3-a_1^3=19.
\]
\checkline{Помню формулу куба суммы и проверяю скобки.}
\[
(x+1)^3=x^3+3x^2+3x+1.
\]
Следовательно,
\[
3a_1^2+3a_1+1=19\quad\Longrightarrow\quad
 a_1^2+a_1-6=0.
\]
\[
(a_1+3)(a_1-2)=0,\qquad a_1=2\quad(a_1>0).
\]
Проверка: $3^3-2^3=27-8=19$.

\checkline{Помню, как меняется объём при изменении масштаба.}
Если каждое ребро увеличить \textbf{в $k$ раз}, то новый объём равен $k^3V$. Например, при увеличении ребра в $3$ раза объём увеличится в $27$ раз.

\sectionhead{Контроль типичных ошибок}
\checkline{В формулу полной поверхности подставляю \textbf{все три} пары измерений и учитываю множитель $2$.}
\checkline{При поиске пространственной диагонали складываю квадраты \textbf{трёх} измерений; у диагонали грани --- двух.}
\checkline{Не путаю $D$ и $D^2$, а также $d$ и $d^2$.}
\checkline{При раскрытии $(a+1)^3$ записываю все четыре слагаемых, включая $+1$.}
\checkline{При решении уравнения о ребре куба исключаю неположительные значения.}
\checkline{Для объёма указываю кубические, для площади --- квадратные единицы.}

\clearpage
\thispagestyle{fancy}
\begin{center}
{\Large\color{darkaccent} Алгоритм: задача на параллелепипед или куб}
\end{center}
\vspace{11mm}
\begin{center}
\begin{tikzpicture}[
  >={Stealth[length=2.7mm]},
  terminal/.style={ellipse,draw=darkaccent,line width=.7pt,minimum width=49mm,minimum height=10mm,align=center,font=\small,inner sep=3pt},
  action/.style={rounded corners=2pt,draw=accent,line width=.65pt,minimum width=87mm,minimum height=14mm,text width=79mm,align=center,font=\small,inner sep=4pt},
  condition/.style={diamond,aspect=2.3,draw=darkaccent,line width=.7pt,text width=47mm,align=center,font=\small,inner sep=1pt},
  arrow/.style={->,line width=.75pt,draw=darkaccent},
  node distance=11mm]
\node[terminal] (start) {Прочитайте вопрос задачи};
\node[action,below=of start] (identify) {Определите фигуру и искомую величину: объём, площадь, диагональ};
\node[action,below=of identify] (formula) {Запишите соответствующую формулу и подставьте известные данные};
\node[condition,below=15mm of formula] (missing) {Не хватает длины ребра?};
\node[action,below=18mm of missing] (solve) {Вычислите искомую величину};
\node[action,right=17mm of missing,text width=39mm,minimum width=47mm] (find) {Найдите ребро через формулу диагонали или условие об изменении объёма};
\node[action,below=of solve] (check) {Проверьте положительность длин, арифметику и единицы измерения};
\node[terminal,below=of check] (end) {Запишите ответ};
\draw[arrow] (start) -- (identify);
\draw[arrow] (identify) -- (formula);
\draw[arrow] (formula) -- (missing);
\draw[arrow] (missing) -- node[fill=white,inner sep=2pt,left,font=\small] {нет} (solve);
\draw[arrow] (missing) -- node[fill=white,inner sep=2pt,above,font=\small] {да} (find);
\draw[arrow] (find.north) |- ([xshift=7mm,yshift=4mm]formula.east) -- (formula.east);
\draw[arrow] (solve) -- (check);
\draw[arrow] (check) -- (end);
\end{tikzpicture}
\end{center}
\vfill
\begin{center}
\color{quiet}\small Сначала определите формулу искомой величины. Если данных недостаточно,\par
восстановите недостающее ребро и вернитесь к формуле.
\end{center}
\clearpage

\begin{center}
{\Large\color{darkaccent} Тренировочный блок}\par
\vspace{2mm}
{\small 10 заданий \textbullet\ от простого к сложному}
\end{center}
\vspace{3mm}
Решите самостоятельно. В задачах с размерами в сантиметрах записывайте единицы площади и объёма. Для всех фигур предполагаются положительные длины рёбер.

\subsection*{Базовый уровень}
\begin{enumerate}[label=\textbf{\arabic*.},itemsep=8pt]
\item Рёбра прямоугольного параллелепипеда равны $3$, $4$ и $5$ см. Найдите его объём.
\item Рёбра прямоугольного параллелепипеда равны $2$, $3$ и $4$ см. Найдите площадь полной поверхности.
\item Объём прямоугольного параллелепипеда равен $72\ \text{см}^3$. Ребро, перпендикулярное некоторой грани, равно $9$ см. Найдите площадь этой грани.
\end{enumerate}
\subsection*{Средний уровень}
\begin{enumerate}[label=\textbf{\arabic*.},start=4,itemsep=8pt]
\item Стороны прямоугольной грани параллелепипеда равны $5$ и $12$ см. Найдите диагональ этой грани.
\item Измерения прямоугольного параллелепипеда равны $3$, $4$ и $12$ см. Найдите его пространственную диагональ.
\item Два ребра прямоугольного параллелепипеда, выходящие из одной вершины, равны $2$ и $6$ см, а его пространственная диагональ равна $7$ см. Найдите объём.
\item Ребро куба равно $4$ см. Найдите его объём и площадь полной поверхности.
\end{enumerate}
\subsection*{Повышенный уровень}
\begin{enumerate}[label=\textbf{\arabic*.},start=8,itemsep=8pt]
\item Пространственная диагональ куба равна $3\sqrt{3}$ см. Найдите площадь полной поверхности куба.
\item Диагональ грани куба равна $5\sqrt{2}$ см. Найдите площадь полной поверхности куба.
\item После увеличения каждого ребра куба на $2$ см его объём увеличился на $98\ \text{см}^3$. Найдите первоначальную длину ребра.
\end{enumerate}
\clearpage

\begin{center}
{\Large\color{darkaccent} Ответы}\par
\vspace{2mm}
{\small Проверьте себя после выполнения тренировочного блока.}
\end{center}
\vspace{10mm}
\begin{center}
\renewcommand{\arraystretch}{1.58}
\begin{tabularx}{.86\linewidth}{@{}>{\centering\arraybackslash}p{.16\linewidth}X@{}}
\toprule
\textbf{№} & \textbf{Ответ}\\
\midrule
1 & $60\ \text{см}^3$\\
2 & $52\ \text{см}^2$\\
3 & $8\ \text{см}^2$\\
4 & $13\ \text{см}$\\
5 & $13\ \text{см}$\\
6 & $36\ \text{см}^3$\\
7 & $64\ \text{см}^3$; $96\ \text{см}^2$\\
8 & $54\ \text{см}^2$\\
9 & $150\ \text{см}^2$\\
10 & $3\ \text{см}$\\
\bottomrule
\end{tabularx}
\end{center}
\vspace{14mm}
\begin{center}
{\large\color{darkaccent} Итоговая памятка}
\end{center}
\begin{center}
\renewcommand{\arraystretch}{1.38}
\begin{tabularx}{\linewidth}{@{}>{\raggedright\arraybackslash}p{.41\linewidth}X@{}}
\toprule
\textbf{Сигнал в условии} & \textbf{Что делать}\\
\midrule
«Объём» & $V=abc$ либо $V=a^3$\\
«Полная поверхность» & $S=2(ab+ac+bc)$ либо $S=6a^2$\\
«Диагональ грани» & Теорема Пифагора для двух рёбер\\
«Диагональ тела» & Теорема Пифагора для трёх рёбер\\
«Ребро увеличили на $k$» & Заменить $a$ на $a+k$; сравнить объёмы\\
«Ребро увеличили в $k$ раз» & Объём увеличился в $k^3$ раз\\
\bottomrule
\end{tabularx}
\end{center}
\vfill
\begin{center}\small\color{quiet}Лёвин Артём Александрович \textbullet\ эксклюзивно для Ксении Клыковой\end{center}
\end{document}
